\chapter{Architectural Decisions}

This chapter records the major architectural decisions made during the evolution of nopCommerce
for the VerdeMart omnichannel scenario. Each decision is traceable to one or more quality
attribute scenarios from Chapter 4 and to the bounded contexts defined in Chapter 3. Rejected
alternatives are included with explicit reasons, since the value of a decision is inseparable
from understanding what was considered and set aside.

\section{ADR 1: Asynchronous Fulfillment Propagation}

\subsection{Context}

After a web order is accepted and paid, nopCommerce must notify the warehouse execution system.
In the current baseline, external communication happens either through synchronous plugin calls
or in-process events, both of which tie the checkout path directly to the availability of the
downstream system.

QAS 1 (Availability under Warehouse Failure) requires that a warehouse timeout or
rejection does not prevent a web order from being recorded. QAS 4 (Recoverability after
Shipping Outage) further requires that failed fulfillment and shipping work is retained and
retried, not silently dropped. A synchronous warehouse call from within the checkout flow
cannot satisfy either constraint.

\subsection{Decision}

Fulfillment propagation is handled asynchronously. When nopCommerce accepts an order, it
records an integration task for the fulfillment request. This task is picked up outside the
customer request path, so a slow or unavailable warehouse has no effect on checkout completion.
If the warehouse rejects or fails to acknowledge the request, the task stays pending and is
retried by the Integration Coordination context.

Shipping follows the same model. Once the warehouse marks an order as ready to dispatch, the
shipping adapter requests the carrier label independently. Provider unavailability results in a
stored, retryable request rather than a visible checkout failure.

This is consistent with the boundary rule established in Chapter 3: checkout does not
synchronously depend on warehouse, POS, or shipping systems.

\subsection{Alternatives Considered}

\textbf{Synchronous API call to the warehouse during checkout.} Calling the warehouse directly
after order creation is straightforward to implement, but it makes checkout latency and success
rate dependent on warehouse response times. During peak traffic, a warehouse response exceeding
ten seconds degrades the buying experience, and an outage blocks all order confirmations
entirely. QAS 1 rules this out.

\textbf{In-process event bus (current nopCommerce approach).} The existing
\texttt{OrderPlacedEvent} fires in-process and awaits registered consumers. It works well for
internal reactions like cache invalidation or plugin hooks, but it offers no durability for
external integration. A process restart between order commit and event dispatch loses the event
with no retry path. Chapter 2 identifies this limitation explicitly as a gap for the target
scenario.

\subsection{Consequences}

Checkout becomes independent of warehouse and shipping availability, which is the primary goal.
The trade-off is that fulfillment state is eventually consistent: the order model must expose
intermediate states such as awaiting warehouse confirmation or fulfillment delayed, and those
states must be visible to support agents and operators. Integration tasks also need idempotency
guarantees to prevent duplicate fulfillment commands when retries occur, which is the subject
of ADR 5.

\section{ADR 2: Transactional Outbox for Durable Integration Task Publishing}

\subsection{Context}

ADR 1 requires that fulfillment work is propagated asynchronously, but that alone does not
solve durability. A common approach publishes a message after the database transaction commits.
The problem is the gap between the two operations: if the process restarts or the broker is
unreachable at that moment, the integration task is lost while the order appears confirmed to
the customer. The warehouse never receives the request, and without explicit detection, nobody
knows.

The Integration Coordination context defined in Chapter 3 owns integration messages,
correlation identifiers, idempotency keys, and retry state. For this ownership to be
meaningful, there must be a reliable way to record integration work at the point the order is
created.

\subsection{Decision}

Integration tasks are written in the same database transaction as the order. When
\texttt{OrderProcessingService} creates an order, it also inserts an outbox record in the same
transaction. The Integration Worker periodically reads unpublished records, dispatches them to
the integration layer, and marks them published. If the process restarts before the worker
runs, the record survives and is processed on the next cycle.

This gives a simple atomicity guarantee: the order and its integration task either both exist
or neither does. At-least-once delivery follows naturally, which means consumers must be
idempotent because duplicate messages can arrive during retries after uncertain acknowledgements.

\subsection{Alternatives Considered}

\textbf{Publish after commit (dual-write).} Committing the order and then publishing to the
broker is the path of least resistance. The gap it creates is not theoretical: a network blip,
broker restart, or process crash at that moment silently loses the integration task. Recovering
requires reconciling the order table against warehouse acknowledgements by hand, which is
fragile and hard to sustain operationally.

\textbf{Scheduled task polling without an outbox table.} A task could scan for orders in a
specific status and attempt delivery without a dedicated outbox record. This sounds simpler but
effectively reconstructs outbox semantics informally. Tracking which orders have been sent and
ensuring the delivery attempt and the status update are atomic requires the same care as an
explicit outbox, but without the clear ownership and structure.

\textbf{Relying on broker-side durability guarantees alone.} Some brokers provide strong
delivery guarantees once a message is accepted. The difficulty is that the producer still needs
to successfully reach the broker within the same logical operation as the order commit. When
the broker is temporarily unavailable at order creation time, the task is never enqueued. An
outbox moves the broker dependency to the Integration Worker, which tolerates broker
unavailability at commit time without affecting the customer.

\subsection{Consequences}

No paid order is silently dropped due to a transient broker or process failure. The worker
polling interval introduces a small propagation delay between order creation and warehouse notification,
which is acceptable given the asynchronous design. The outbox table becomes
integration-critical infrastructure and must be covered by database backups and operational
monitoring. The publishing interval also determines the minimum propagation latency, which must
remain within the bounds set by QAS 1.

\section{ADR 3: Adapter Pattern with No Shared Database Across External Boundaries}

\subsection{Context}

VerdeMart's operational ecosystem includes at minimum a warehouse execution system and a
shipping provider, with POS and ERP systems also present. As the business grows, further
operational systems may be added or replaced. Each one has its own protocols, data models, and
failure behaviour.

The nopCommerce database is a single schema that mixes catalog, orders, customers, inventory,
shipments, configuration, and logs. Allowing external systems to query or write this schema
directly would create implicit dependencies on internal table structures, making it difficult
to evolve either side independently. It also creates ambiguous data ownership: if the warehouse
system writes directly to the nopCommerce stock table, it is unclear which system is the
authoritative source and which is a consumer.

QAS 5 (Modifiability for New Operational Channels) requires that a new system can be
integrated through no more than one adapter boundary, without touching checkout logic or
coupling to the nopCommerce schema.

\subsection{Decision}

Each external system is connected through a dedicated adapter that translates between the
internal integration contract and the system's specific protocol, data model, and error
conventions. The Warehouse, Inventory, Shipping, and Store Operations adapters are each
independently replaceable. State flowing back into nopCommerce, fulfillment status, stock
updates, carrier progress, arrives through events or adapter APIs and is stored as a
projection with source and freshness metadata. No external system reads or writes the
nopCommerce database directly.

\subsection{Alternatives Considered}

\textbf{Direct database access for external systems.} Giving the warehouse system or ERP read
or write access to the nopCommerce database is the lowest-effort path for a small and stable
set of integrations. The problem is that it couples every external system to the internal
schema. Any table rename, column change, or migration that makes sense for the commerce core
becomes a cross-team coordination problem. It also makes data ownership opaque, which
complicates debugging and reconciliation when systems disagree.

\textbf{A single generic integration service.} Routing all external events through one
integration service reduces the number of deployable components. The cost is that a single
service must handle the distinct failure modes, retry policies, and domain translations for
every external system simultaneously. A shipping provider outage and a warehouse delay are
different problems requiring different handling; merging them into one service makes both
harder to reason about and monitor independently.

\textbf{Plugin calls from within the nopCommerce process.} nopCommerce plugins can call
external systems, and this pattern is already used for payment and shipping-rate providers.
For the omnichannel integration boundaries in Scenario C, however, plugins inherit the same
limitations as synchronous calls: they block the request path, they have no built-in retry or
dead-letter mechanism, and their availability is tied to the monolith process. They remain the
right tool for internal extensibility, not for durable external coordination.

\subsection{Consequences}

Each external system has a bounded, explicit integration contract. Replacing a warehouse
provider means replacing its adapter; the commerce core and the other adapters are unaffected.
New channels follow the same pattern: one new adapter, one set of event subscriptions, no
changes to checkout. The adapter boundary also enforces the language separation noted in
Chapter 3: warehouse, POS, and carrier terminology does not leak into the core order model.
The cost is that every new system needs an adapter implementation, but the scope of that work
is predictable and isolated.

\section{ADR 4: nopCommerce Monolith Retained as the Commerce Core}

\subsection{Context}

When an omnichannel scenario adds warehouse, POS, shipping, and ERP integrations, a natural
design instinct is to decompose the central platform into smaller services. nopCommerce already
separates concerns into layers and the plugin model gives it modularity at the code level, so
the question of whether to extract services from it is worth taking seriously.

The quality attribute scenarios in Chapter 4 identify five pressures: availability,
consistency, performance, recoverability, and modifiability. Looking at each one, none of them
requires extracting checkout, order management, inventory, or catalog from the monolith. The
existing layered architecture already provides clear seams. Breaking those capabilities into
separate services would introduce distributed transaction problems and network dependencies that
do not currently exist, without resolving any of the five identified pressures.

\subsection{Decision}

nopCommerce remains the commerce core and is not decomposed. Checkout, catalog, payments,
customer accounts, and canonical order management stay inside the monolith. The main new
independently deployable platform component introduced by this evolution is the Integration
Coordination context, which owns the outbox, retry, idempotency, and reconciliation behaviour
described in ADR 2. Adapters for warehouse, shipping, inventory, and store operations sit around this layer.
The monolith interacts with the Integration Coordination context through its outbox table and a
publishing worker, remaining decoupled from the internal design of any individual
adapter.

\subsection{Alternatives Considered}

\textbf{Extract order management as an independent service.} Separating order creation, payment
result handling, and order history into a standalone service would produce an independent
deployment unit, but it removes a large body of stable, well-tested functionality from the
monolith. More importantly, it would introduce a synchronous dependency between the storefront
and the new order service at the point of checkout, which directly conflicts with the
availability requirement in QAS 1. The decomposition adds complexity without improving
any of the measured quality attributes.

\textbf{Extract inventory as an independent service.} Given QAS 2 (Consistency under
Stale Inventory), a dedicated inventory service was considered. The scenario requires
cross-channel stock visibility with freshness metadata and stale-state detection. The Inventory
Adapter already provides this by updating projections inside nopCommerce through integration
messages, extending the existing stock model rather than replacing it. Extracting inventory
into a separate service would require network calls for every stock check at checkout, a new
database to own the projection, and a service boundary that the adapter model renders
unnecessary.

\textbf{Full microservices rewrite.} Decomposing nopCommerce into catalog, cart, checkout,
payment, order, inventory, and notification services would require recreating a large volume of
working functionality across independently operated services, coordinating distributed
transactions across service boundaries, and shifting most of the project effort from
architectural evolution to infrastructure. The result would be a more complex system with the
same functional scope.

\subsection{Consequences}

The monolith retains its operational simplicity: one deployable unit for the full
customer-facing commerce experience, with all existing capabilities, checkout, payments,
catalog, admin, and plugin extensibility intact. The Integration Coordination context is tightly
bounded to reliability and coordination, while adapter services remain boundary-specific and
replaceable. The main limitation is that the commerce core cannot be scaled independently from
its own checkout and order-management data model, but this is sufficient for the load levels
identified in QAS 3 and is not a constraint that justifies decomposition at this stage.
This decision does not preclude future extraction of read-heavy services such as the catalog or
order-status projection if load patterns justify it later.

\section{ADR 5: Idempotency Key Strategy for Integration Messages}

\subsection{Context}

ADR 1 establishes asynchronous fulfillment, and ADR 2 guarantees at-least-once delivery
through the outbox. At-least-once delivery means that the same fulfillment or label-creation
message may be delivered more than once during retries. If the warehouse or shipping adapter
processes the same message twice, the result can be a duplicate fulfillment request or a
second shipping label for the same order, both of which cause operational problems.

\subsection{Decision}

Every integration message carries an \texttt{IdempotencyKey} generated at message creation
time as a UUID. The key is stored in the outbox record so it is stable across retries. Each
adapter is responsible for detecting a duplicate key and returning a success acknowledgement
without repeating the side-effecting operation. The Integration Coordination context records
key-to-outcome pairs for a configurable deduplication window (default: 48 hours). Successful
outcomes are cached to prevent unnecessary processing.

\subsection{Alternatives Considered}

\textbf{Order-ID-based idempotency.} Using the nopCommerce order identifier as the idempotency
key is simpler, but it conflates the order identity with the delivery attempt identity. If an
order requires two separate fulfillment requests for different warehouse splits, both would
carry the same key and the second would be incorrectly treated as a duplicate.

\textbf{Idempotency enforced only at the database level.} Some adapters add a unique constraint
on the order identifier in their local database. This prevents duplicate processing but ties
the idempotency guarantee to adapter internals rather than the message contract. It cannot be
relied upon for external systems where the adapter is a surrogate or stub.

\subsection{Consequences}

Adapters must implement idempotency checks on the \texttt{IdempotencyKey}. The deduplication
window creates a bounded storage requirement for the key-outcome table. The risk of missing a
duplicate outside the 48-hour window is low for the use case, since retries for any single
message should complete well within that period.

\section{ADR 6: Retry Policy and Circuit Breaker for Adapter Calls}

\subsection{Context}

When a warehouse or shipping adapter is unavailable, the Integration Worker retries the failed
message. Without a controlled retry policy, the worker can overwhelm a recovering adapter with
a burst of accumulated retries, turning a temporary outage into a sustained overload. A circuit
breaker pattern can prevent this by stopping retry attempts after a threshold is reached and
resuming only after a health check succeeds.

\subsection{Decision}

The retry policy applies exponential backoff with the following parameters: initial delay of
2 seconds, backoff multiplier of 2, maximum delay of 5 minutes, and maximum of 10 retry
attempts. After 5 consecutive failures from the same adapter, the circuit breaker opens and
blocks further attempts for a configurable cooldown period (default: 5 minutes). After the
cooldown, one probe attempt is made. If it succeeds, the circuit closes and normal retry
resumes. If it fails, the cooldown restarts.

Messages that exhaust all retry attempts are moved to the dead-letter queue described in
ADR 7.

\subsection{Alternatives Considered}

\textbf{Fixed-interval retry without a circuit breaker.} Fixed-interval retries are simpler but
can create a thundering-herd effect when a dependency recovers. They also provide no protection
against sustained outages: a dependency down for 30 minutes would receive hundreds of retry
attempts per message before the outbox is drained.

\textbf{Linear backoff.} Linear backoff reduces the thundering herd compared to fixed intervals
but converges to a high retry rate more quickly than exponential backoff and does not provide
the same degree of protection during long outages.

\subsection{Consequences}

The maximum propagation delay for a fulfilled order under sustained warehouse failure is bounded
by the backoff ceiling and max retries (approximately 50 minutes of active retry before
dead-letter). The circuit breaker protects the recovering adapter from burst load. Operators
must be informed when a circuit opens, because it means fulfillment work has stopped flowing
to the affected adapter. This requires a monitoring alert and a visible circuit-state indicator
in the Support and Operations View.

\section{ADR 7: Dead-Letter Queue and Operator Intervention Model}

\subsection{Context}

After the maximum retry attempts defined in ADR 6 are exhausted, the integration message has
failed permanently in the automated path. A decision must be made about what happens next:
the message can be silently discarded, moved to an archive, or presented to an operator for
manual action.

Discarding messages is not acceptable for fulfillment or shipping work, since a lost
fulfillment message means a paid order is never sent to the warehouse. The architecture
therefore needs a durable dead-letter record and a defined escalation process.

\subsection{Decision}

Failed messages are written to a dead-letter table in the nopCommerce database with the
following fields: the original message payload, the idempotency key, the correlation identifier
(order ID), the adapter that failed, the failure reason for the last attempt, the timestamp of
the first and last attempt, and the escalation state (new, acknowledged, resolved).

The Support and Operations View surfaces dead-letter records alongside pending and retrying
work. An alert fires to the operations team when a new dead-letter entry is created. Operators
can manually re-queue a dead-letter message after the dependency is confirmed healthy, or mark
it as resolved if the situation is handled by other means such as manual fulfillment or refund.

Dead-letter records are retained for 30 days and then archived.

\subsection{Alternatives Considered}

\textbf{Separate message broker dead-letter queue.} RabbitMQ supports dead-letter exchanges
natively, and routing failed messages there is a common pattern. For this project, keeping the
dead-letter record in the database gives operators a single interface for all integration work
rather than requiring access to the broker management UI. It also makes dead-letter records
part of the same backup and audit trail as orders.

\textbf{Automatic retry without a dead-letter limit.} Retrying indefinitely avoids requiring
operator intervention but creates unbounded queue growth and makes it impossible to distinguish
transient from permanent failures. Operations cannot act on a problem they cannot see.

\subsection{Consequences}

The dead-letter table is integration-critical infrastructure and must be monitored and backed
up alongside the orders table. The operator intervention model introduces a human step in the
recovery path, which is appropriate for permanent failures but requires clear UI affordances
in the Support and Operations View. The 30-day retention policy must be aligned with the
business's order dispute window.
